%! Simplifying Rational Expressions — Unit 1, Lesson 1.7 — Homework
\documentclass[10pt]{article}
\usepackage{atda-article}
\usepackage{atda-boxes}
\newcommand{\TallMath}[1]{$\displaystyle #1\rule[-1.4em]{0pt}{3.2em}$}
\setlength{\workrowsep}{6pt}

\begin{document}

\pageheader{Unit 1, Lesson 1.7}{Homework}
% NAMESTRIP: no \namedateperiod here. The student writes their name once, on the
% cover this component is stapled behind. See references/conventions.md.

\vspace{0.15in}

\boxguard[16]
\begin{notesbox}{Practice --- Factor First, Then Cancel Factors Only}
\small
Show your work. Every answer needs two parts: the \textbf{simplified expression} and \textbf{every
excluded value}, read from the denominators you started with.

\begin{enumerate}[leftmargin=1.6em, itemsep=6pt, topsep=4pt]

\item State every excluded value: \ \TallMath{\frac{x-1}{x+7}}
\begin{work}
  x+7 &= 0 \quad\Rightarrow\quad x = -7
\end{work}

\item State every excluded value: \ \TallMath{\frac{x+4}{x^2-36}}
\begin{work}
  (x+6)(x-6) &= 0 \\
  x &= -6 \quad \text{or} \quad x = 6
\end{work}

\item Simplify: \ \TallMath{\frac{8x^3}{12x}}
\begin{work}
  \frac{8x^3}{12x} &= \frac{4x \cdot 2x^2}{4x \cdot 3} = \frac{2x^2}{3} \qquad (x \ne 0)
\end{work}

\item Simplify: \ \TallMath{\frac{x^2-49}{x^2-14x+49}}
\begin{work}
  \frac{x^2-49}{x^2-14x+49} &= \frac{(x+7)(x-7)}{(x-7)(x-7)} \\
      &= \frac{x+7}{x-7} \qquad (x \ne 7)
\end{work}

\item Simplify: \ \TallMath{\frac{x^2+6x+8}{x^2-4}}
\begin{work}
  \frac{x^2+6x+8}{x^2-4} &= \frac{(x+2)(x+4)}{(x+2)(x-2)} \\
      &= \frac{x+4}{x-2} \qquad (x \ne -2, \; x \ne 2)
\end{work}

\item Simplify: \ \TallMath{\frac{5x-20}{x^2-16}}
\begin{work}
  \frac{5x-20}{x^2-16} &= \frac{5(x-4)}{(x+4)(x-4)} \\
      &= \frac{5}{x+4} \qquad (x \ne -4, \; x \ne 4)
\end{work}

\item Multiply and simplify: \ \TallMath{\frac{x+5}{x-3} \cdot \frac{x^2-9}{x^2-25}}
\begin{work}
  \frac{x+5}{x-3} \cdot \frac{(x+3)(x-3)}{(x+5)(x-5)}
      &= \frac{x+3}{x-5} \qquad (x \ne 3, \; x \ne 5, \; x \ne -5)
\end{work}

\item Divide and simplify: \ \TallMath{\frac{x^2-4}{6x} \div \frac{x+2}{3x^2}}
\begin{work}
  \frac{(x+2)(x-2)}{6x} \cdot \frac{3x^2}{x+2}
      &= \frac{x(x-2)}{2} \qquad (x \ne 0, \; x \ne -2)
\end{work}

\item Subtract and simplify: \ \TallMath{\frac{x^2}{x+6} - \frac{36}{x+6}}
\begin{work}
  \frac{x^2-36}{x+6} &= \frac{(x+6)(x-6)}{x+6} \\
      &= x-6 \qquad (x \ne -6)
\end{work}

\end{enumerate}
\end{notesbox}

\vspace{0.10in}

\boxguard[16]
\begin{scenariobox}[In Context --- The Community Garden Bed]{royal}
\small
A rectangular raised bed at the community garden covers $A(x) = 2x^2+14x+20$ square feet, and the
plan sheet lists its \textbf{width} as $(x+2)$ feet.

\begin{enumerate}[leftmargin=1.6em, itemsep=6pt, topsep=4pt, start=10]

\item Factor $A(x)$ completely. (GCF first.)
\begin{work}
  A(x) &= 2x^2+14x+20 \\
       &= 2\left(x^2+7x+10\right) = 2(x+5)(x+2)
\end{work}

\item Find the bed's \textbf{length} by simplifying \ \TallMath{\frac{A(x)}{x+2}}. State the
excluded value, and say why the garden plan never reaches it.
\begin{work}
  \frac{A(x)}{x+2} &= \frac{2(x+5)(x+2)}{x+2} = 2x+10 \quad \text{feet} \\
  x+2 &= 0 \; \Rightarrow \; x = -2 \; \text{excluded --- but a width must be positive}
\end{work}

\item The plan uses $x=3$. Give the bed's width, length, and area, and check against $A(3)$. Then
the mulch bill is $M(x) = 6x^2+42x+60$ dollars: simplify \ \TallMath{\frac{M(x)}{A(x)}} \ and say
in one sentence what your answer means to the gardeners.
\begin{work}
  \text{width} &= 5 \text{ ft}, \quad \text{length} = 16 \text{ ft}, \quad 5 \cdot 16 = 80 = A(3) \quad\checkmark \\
  \frac{M(x)}{A(x)} &= \frac{3\left(2x^2+14x+20\right)}{2x^2+14x+20} = 3,
      \qquad \text{at } x=3: \; \frac{240}{80} = 3 \quad\checkmark
\end{work}
\writeline

\end{enumerate}
\end{scenariobox}

\vspace{0.10in}

\boxguard[16]
\begin{scenariobox}[A First Look Ahead]{goldacc}
\small
\begin{enumerate}[leftmargin=1.6em, itemsep=6pt, topsep=4pt, start=13]

\item Unit 2 studies \emph{functions}. Consider \ \TallMath{f(x) = \frac{x+1}{x-3}}.
\begin{enumerate}[label=(\alph*), leftmargin=1.6em, itemsep=4pt, topsep=3pt]
  \item State the excluded value of $f$.
\begin{work}
  x-3 &= 0 \quad\Rightarrow\quad x = 3
\end{work}
  \item The \textbf{domain} of a function is the set of inputs it accepts. In one sentence, say
        what your answer to (a) tells you about the domain of $f$ --- and what you would expect the
        graph of $f$ to do near that value.

\writeline
\end{enumerate}

\end{enumerate}
\end{scenariobox}

\vspace{0.10in}

\boxguard[18]
\begin{extensionbox}
\small
\textbf{Unlike denominators, when one factors into the other.} Simplify
\ \TallMath{\frac{3}{x-2} + \frac{5}{x^2-4}} \ into a single fraction and state every excluded
value. Then \emph{verify} your answer by evaluating both the original sum and your result at $x=3$.
\begin{work}
  \frac{3}{x-2} + \frac{5}{(x+2)(x-2)}
      &= \frac{3(x+2)}{(x+2)(x-2)} + \frac{5}{(x+2)(x-2)} \\
      &= \frac{3x+11}{(x+2)(x-2)} \qquad (x \ne 2, \; x \ne -2) \\
  \text{at } x=3: \; \frac{3}{1} + \frac{5}{5} = 4, \quad &\frac{20}{(5)(1)} = 4 \quad\checkmark
\end{work}
\writeline
\end{extensionbox}

\vspace{0.10in}

\begin{remindbox}
\small
\textbf{That closes Unit 1.} Every lesson since 1.3 has been the same sentence in a different
costume: \emph{factoring turns a sum into a product, and a product is something you can act on.} A
product tells you the dimensions (1.3--1.5), it tells you when a height is zero (1.6), and today it
tells you what cancels and what is forbidden. Review the whole unit for the test --- and keep the
habit that ran through all of it: \textbf{check whether your answer is complete}, not just true.
\end{remindbox}

\end{document}
